\documentclass[11pt]{article}
\usepackage{varnernotes}

\begin{document}

\lecturenotes{11b}{Protective Collars with Standard Listed Options---TJUL as a Comparison}{Fall 2026}{November 5, 2026}{Jeffrey D. Varner}

\learningobjectives
  {Construct a standard listed collar}
  {Match 100 shares with one ordinary listed put and one ordinary listed call using a common expiration.}
  {Compute economic risk and reward}
  {Use executable premiums to calculate the floor, cap, breakeven, maximum loss, and maximum profit.}
  {Compare implementation choices}
  {Distinguish a self-managed standard collar from TJUL's FLEX-option fund structure without requiring students to trade FLEX options.}

\section{The Core Position Uses Standard Listed Options}

A protective collar combines long exposure to an underlying asset, a long put
that limits downside, and a short call that helps finance the put by surrendering
upside. The core classroom construction uses ordinary listed options that an
options-approved brokerage account may support. FLEX options are not required.

Let one standard option contract represent $m=100$ shares and hold $n$ matched
contracts. At time zero, buy or already own $mn$ shares at $S_0$, buy $n$ puts
with strike $K_P$, and sell $n$ calls with strike $K_C$, where both option legs
have the same expiration and
\begin{equation}
  \label{eq:strike-order}
  K_P<S_0<K_C.
\end{equation}
The terminal value is
\begin{equation}
  \label{eq:collar-value}
  V_T=mn\left[S_T+(K_P-S_T)^+-(S_T-K_C)^+\right].
\end{equation}
Therefore,
\begin{equation}
  \label{eq:collar-piecewise}
  V_T=mn
  \begin{cases}
    K_P, & S_T<K_P,\\
    S_T, & K_P\leq S_T\leq K_C,\\
    K_C, & S_T>K_C.
  \end{cases}
\end{equation}
The put creates the floor, the shares participate one-for-one between the
strikes, and the short call creates the cap.

\begin{figure}[ht]
  \centering
  \begin{tikzpicture}[x=0.78cm,y=0.62cm,>=Latex,font=\sffamily\small]
    \draw[->,vnmuted] (0,0)--(12.4,0) node[right] {$S_T$};
    \draw[->,vnmuted] (0,0)--(0,7.1) node[above] {$V_T$};
    \draw[vnlink,very thick] (0.4,2.1)--(3.5,2.1)--(9.0,6.0)--(11.8,6.0);
    \draw[vncarnelian,dashed] (3.5,0)--(3.5,2.1) node[below] {$K_P$};
    \draw[vncarnelian,dashed] (9.0,0)--(9.0,6.0) node[below] {$K_C$};
    \node[vnmuted,align=center] at (1.8,3.0) {put floor};
    \node[vnmuted,align=center] at (6.3,4.8) {share\\participation};
    \node[vnmuted,align=center] at (10.4,6.7) {call cap};
  \end{tikzpicture}
  \caption{Terminal value of a standard listed protective collar. Premiums
  shift the profit diagram but do not change the three payoff regions.}
  \label{fig:collar-payoff}
\end{figure}

The Options Clearing Corporation reports that each ordinary standard equity
option generally covers 100 shares and that standard equity options are
American-style. Corporate actions can create adjusted contracts with a
different deliverable, so the broker's contract specification must still be
checked.

\section{Build the Retail Order Ticket}

For an existing 100-share position, the opening option order is:
\begin{enumerate}[leftmargin=1.5em,itemsep=2pt,topsep=3pt]
  \item buy to open one put at strike $K_P$; and
  \item sell to open one call at strike $K_C$ with the same expiration.
\end{enumerate}
For a new position, buy 100 shares as well. Some brokers support a three-leg
stock-plus-options combination order. Otherwise, entering the legs separately
creates execution and interim-exposure risk.

Use executable quote sides in the pre-trade estimate. Let $P_0^{ask}$ be the
put ask and $C_0^{bid}$ the call bid. A conservative initial cost for one
matched collar is
\begin{equation}
  \label{eq:initial-cost}
  \Pi_0=m\left(S_0+P_0^{ask}-C_0^{bid}\right),
  \qquad m=100.
\end{equation}
The broker may fill a multi-leg limit order inside the component quotes, but a
midpoint fill is an execution outcome rather than an assumption. Verify the
underlying, deliverable, strikes, expiration, multiplier, option approval,
buying power, and ex-dividend date before entering the order.

\notebooklink
  {L11b: Standard-Listed SPY Protective Collar}
  {https://github.com/varnerlab/CHEME-5660-CourseRepository-Fall-2026/blob/main/lectures/week-11/L11b/CHEME-5660-L11b-Example-StandardListed-SPY-ProtectiveCollar-Fall-2026.ipynb}
  {Match 100 illustrative SPY shares with one listed put and call, use executable quote sides, calculate risk metrics, and compare the result with TJUL.}

\section{Payoff Is Not Profit}

Ignoring financing, dividends, taxes, and transaction fees, expiration profit
is $V_T-\Pi_0$. For a matched collar,
\begin{align}
  \label{eq:risk-metrics}
  \text{maximum loss}
    &=\Pi_0-mnK_P,\\
  \text{maximum profit}
    &=mnK_C-\Pi_0,\\
  \text{interior breakeven}
    &=\frac{\Pi_0}{mn}
      =S_0+P_0^{ask}-C_0^{bid}.
\end{align}
The breakeven formula applies when it lies between the strikes. A ``zero-cost
collar'' only means that the quoted put premium is approximately financed by
the quoted call premium. It does not remove underlying risk, bid--ask cost,
assignment risk, taxes, financing, or foregone upside.

\section{The Standard Chain Constrains the Design}

Standard listed options restrict the investor to available strikes and
expirations. Raising $K_P$ improves the floor but generally increases the put
cost. Lowering $K_C$ generally collects more call premium but surrenders upside
sooner. A reproducible selection rule can maximize $K_C$ subject to a minimum
$K_P$, maximum net debit, minimum volume/open interest, and maximum spread.

For example, a 10\%-down/10\%-up collar has
\begin{equation}
  \label{eq:ten-ten}
  K_P=0.90S_0,
  \qquad K_C=1.10S_0.
\end{equation}
It does not provide a 100\% downside buffer. The position participates in the
first 10\% decline and is protected only below $K_P$. A put nearer $S_0$
provides stronger protection but normally requires a larger debit or a lower
call cap.

\section{Keep Growth and Price Changes Distinct}

For a horizon $\Delta t$ measured in years, this course uses
\begin{equation}
  \label{eq:growth-convention}
  r_{t,\Delta t}
  :=\log\!\left(\frac{S_{t+\Delta t}}{S_t}\right)
  =(\Delta t)g_{t,\Delta t}.
\end{equation}
Option diagrams and defined-outcome documents generally quote a cumulative
simple price change. Denote it by a capital $R$:
\begin{equation}
  \label{eq:simple-change}
  R_{t,\Delta t}
  :=\frac{S_{t+\Delta t}-S_t}{S_t}
  =\exp\!\left((\Delta t)g_{t,\Delta t}\right)-1.
\end{equation}
A 90-day collar outcome or a two-year ETF cap is cumulative over its stated
horizon, not an annual growth rate.

\section{Expiration and Early Assignment}

Below $K_P$, put exercise or assignment can sell the matched shares at the floor
strike. Between the strikes, both options normally expire without exercise and
the investor retains the shares. Above $K_C$, call assignment sells the shares
at the cap strike.

Standard equity and ETF options are generally American-style, so assignment can
occur before expiration. Short-call assignment risk is especially relevant near
an ex-dividend date when remaining extrinsic value is small. The operator should
decide before entry how to handle dividends, closing, rolling, exercise,
assignment, and any mismatch between shares and adjusted deliverables. Adjusting
one leg must not accidentally leave an unintended naked option.

\section{TJUL Is the Comparison, Not the Required Instrument}

The Innovator Equity Defined Protection ETF (TJUL) packages a defined outcome
inside an exchange-traded fund. TJUL uses customized FLEX options on SPY; the
disclosed construction uses purchased call exposure, a purchased put, and a
sold call that creates the cap. FLEX options permit customized strikes,
expirations, and exercise styles and use specialized exchange execution
workflows. Students do not need that access. They can:
\begin{itemize}[leftmargin=1.3em,itemsep=2pt,topsep=2pt]
  \item build a self-managed collar with ordinary listed SPY options;
  \item buy or analyze TJUL shares as an ETF product; or
  \item compare the two without trying to replicate TJUL's FLEX positions.
\end{itemize}

Let $R_{SPY}$ be the complete-period cumulative SPY price change and $c$ the
starting cap. A simplified full-buffer outcome before expenses is
\begin{equation}
  \label{eq:defined-outcome}
  R_{TJUL}^{\mathrm{gross}}=
  \begin{cases}
    0, & R_{SPY}<0,\\
    R_{SPY}, & 0\leq R_{SPY}\leq c,\\
    c, & R_{SPY}>c.
  \end{cases}
\end{equation}
This is a complete-period outcome map, not a daily guarantee, a mid-period NAV
model, or an instruction to trade FLEX contracts.

\section{A Dated TJUL Snapshot}

The official materials checked on August 4, 2026 describe an outcome period
from July 1, 2025 through June 30, 2027, a 13.61\% starting cap before expenses,
a 100\% starting buffer before expenses, and a 0.79\% annual management fee.
The prospectus illustrates a 12.03\% cap and 98.42\% buffer after two years of
management fees. Refresh the
\href{https://www.innovatoretfs.com/etf/?ticker=TJUL}{official fund page},
\href{https://www.innovatoretfs.com/pdf/tjul_prospectus.pdf}{prospectus}, and
\href{https://www.innovatoretfs.com/pdf/tjul_factsheet.pdf}{factsheet} before
class. Starting outcomes apply only from the first through the last day of the
outcome period.

\begin{center}
\small
\renewcommand{\arraystretch}{1.18}
\begin{tabular}{p{0.22\textwidth}p{0.34\textwidth}p{0.34\textwidth}}
\hline
\textbf{Feature} & \textbf{Self-managed standard collar} & \textbf{TJUL}\\
\hline
Investor accesses & Shares plus listed put and call & ETF shares\\
Options used & Standard listed contracts & FLEX contracts managed inside fund\\
Granularity & Usually 100-share blocks & Whole or fractional ETF shares, broker dependent\\
Floor and cap & Available strikes, premium, expiration & Dated starting buffer and cap\\
Management & Investor handles execution and assignment & Fund manages option package\\
Costs & Spreads, fees, financing, taxes & Expense ratio, ETF spread, premium/discount, taxes\\
\hline
\end{tabular}
\end{center}

TJUL references the SPY \emph{price} return. Direct SPY ownership may also
receive distributions, so compare complete wealth using consistent price-return
or total-return conventions.

\section{Optional Advanced Track}

An advanced standard-options project can load a dated SPY chain and optimize
collar strikes under floor, cap, liquidity, and maximum-debit constraints. The
previous Wheel lecture and supporting examples remain in the \texttt{advanced/}
subfolder as a separate optional extension.

\Needspace{0.34\textheight}
\keytakeaways
  {In this lecture, we constructed a retail-replicable collar using ordinary listed options and used TJUL only as a defined-outcome comparison.}
  {FLEX options are not required}
  {The core position is 100 shares, one standard listed put, and one standard listed call with a common expiration.}
  {Use executable premiums}
  {Put ask minus call bid determines a conservative initial debit and therefore shifts breakeven, maximum loss, and maximum profit.}
  {Implementation determines the outcome}
  {Listed strikes constrain a self-managed floor and cap; TJUL instead manages FLEX options internally while the investor accesses ETF shares.}
  {Next, generalize sequential decisions to exploration--exploitation problems with stochastic multi-armed bandits.}

\vfill
{\sffamily\footnotesize\color{vnmuted}\textbf{Educational use and risk notice.}
These notes are for instruction only and do not constitute investment advice or
an order recommendation. Options involve risk and require broker approval.}

\end{document}
